\subsection{Real-robot fine-tuning on stair-based manipulation tasks}
\label{sec:exp-real-robot}

We extend the study to physical manipulation on an OpenArm platform, asking
whether post-training can improve a demonstration-trained policy when the
feedback available after each attempt is a human success/failure label.
The hardware pipeline and demonstrations are available, while the controlled
four-way success-rate comparison remains in progress. The episode counts in
\cref{tab:real-robot-budgets} specify the \emph{planned study budgets}, not a
retrospective claim about the data used by exploratory checkpoints. We will
verify checkpoint provenance and retained-episode counts against these budgets
before reporting the comparison. This study tests refinement of an existing
skill rather than learning a new policy from scratch.

\paragraph{Robot hardware and sensing.}
We use an OpenArm leader--follower platform, operating only the left follower
arm for autonomous task execution. This arm has seven actuated joints and a
parallel gripper; the right follower arm remains at its initialized pose.
An operator uses the corresponding leader arm to collect teleoperated
demonstrations. A chest-mounted RealSense D435 provides a fixed overview of
the staircase and target, while a RealSense D405 mounted at the left wrist
provides a close view of the gripper--object interaction. Both cameras deliver
$640\!\times\!480$ RGB images at a nominal 15\,Hz, resized to
$128\!\times\!128$ for the policy. Observations also include the seven left-arm
joint positions, gripper position, and a fixed task-language embedding.
Autonomous inference uses follower proprioception; the stationary leader is
not used as its state estimate. Finger tactile measurements can be recorded
for analysis, but are not inputs to the policies in this comparison.

Inference runs on a local workstation with an NVIDIA RTX 3080 Ti (12\,GB),
using PyTorch inside a Docker-based ROS~2 Humble environment. All four
conditions share this robot, sensor interface, and low-level position control.

\paragraph{Task setup.}
We consider two tasks:
\emph{Jigglypuff into the case}, in which the robot grasps the toy from a stair
tread and releases it into an adjacent open case; and \emph{red-on-purple
stacking}, in which it picks up the red cube and places it on the purple cube.
Both use a hand-made staircase. The stepped support surfaces introduce
different approach heights and nearby edges, making grasping and placement
sensitive to alignment; stacking additionally requires the upper cube to be
positioned accurately before release. \Cref{fig:real-robot-tasks} illustrates
the tasks through successful teleoperated demonstrations from both views.

\paragraph{Episode budgets.}
For Jigglypuff placement, the target base-training budget is 40 teleoperated
demonstration episodes. For red-on-purple stacking, we allocate 80 because
the smaller placement target calls for finer alignment. In each task,
BC refinement uses 40 \emph{additional successful training trajectories},
separate from the base demonstrations. The two RL methods instead share a
pool of 40 physical robot rollout attempts, retaining both successes and
failures and using the same outcome labels. These counts refer to complete
episodes, not optimization epochs, minibatches, or evaluation trials.

Forty retained successful BC trajectories need not correspond to forty robot
attempts: if they are selected from autonomous rollouts, the attempts required
to obtain them will also be reported. Consequently, we do not claim matched
collection effort between the BC and RL refinement budgets. The controlled
actor-update comparison is DICE-RL versus DICE-RL $+$ \method{}, which shares
the same 40-attempt pool and demonstration-replay settings.

\begin{table}[t]
  \centering
  \caption{\textbf{Target training-data budgets for the real-robot study.}
  Counts are episodes. BC refinement starts from the corresponding base and
  uses additional positive trajectories; the two RL methods use the same
  mixed-outcome rollout pool. These are planned budgets to be verified before
  reporting results, not counts of optimization epochs or held-out test trials.}
  \label{tab:real-robot-budgets}
  \small
  \setlength{\tabcolsep}{5pt}
  \begin{tabular}{lccc}
    \toprule
    Task & \makecell{Base BC\\demonstrations} &
      \makecell{Additional positive\\BC trajectories} &
      \makecell{Shared RL\\rollout attempts} \\
    \midrule
    Jigglypuff into case & 40 & 40 & 40 \\
    Red-on-purple stacking & 80 & 40 & 40 \\
    \bottomrule
  \end{tabular}
\end{table}

\begin{figure}[t]
  \centering
  \includegraphics[width=\linewidth]{figures/real_robot_tasks.pdf}
  \caption{\textbf{Two real-robot tasks on a hand-made staircase.}
  Each task is shown as a five-stage sequence, with chest and wrist views in
  paired rows: initial arrangement, grasp, transfer, alignment, and release.
  The wrist view shows the toy released into the case in (a) and the red cube
  resting on the purple cube with the gripper open in (b).
  Frames are extracted from successful \emph{teleoperated demonstrations},
  not autonomous evaluations of the compared methods. Times are nominal video
  times at 15\,fps; the two views are approximately aligned by frame index.
  A consistent top-background crop is applied to all panels.}
  \label{fig:real-robot-tasks}
\end{figure}

\paragraph{Pretrained policy and refinement methods.}
For each task, we first train a flow-matching BC policy on teleoperated
demonstrations. The current policy predicts 30-step chunks of absolute
left-arm joint and gripper targets at 15\,Hz, using a 10-step Euler sampler.
Each control is eight-dimensional: seven joint-position targets and one
gripper-position target. The image encoder is a shared ResNet-18; visual,
proprioceptive, and language features condition the DiT velocity model.
The comparison uses one fixed base checkpoint per task and four conditions:
\begin{itemize}
  \item \textbf{Base FM:} deploy the demonstration-trained checkpoint without
  online refinement.
  \item \textbf{Success-only BC refinement:} initialize from that same
  checkpoint and continue supervised flow-matching training using only
  the 40 additional successful trajectories in
  \cref{tab:real-robot-budgets}. Failed attempts are excluded from these
  imitation targets.
  \item \textbf{DICE-RL:} fine-tune a residual actor and critic over the same
  frozen base using both successful and failed online rollouts, with the
  standard DICE-RL actor objective~\citep{dicerl2026}.
  \item \textbf{DICE-RL $+$ \method{}:} use the same residual RL setup and both
  rollout outcomes, replacing the actor update with \method{}'s clipped,
  base-restoring action-target regression.
\end{itemize}
Thus all three refinement conditions start from pretrained weights; none
trains a task policy from scratch. Success-only BC tests how much improvement
can be obtained by imitating successful practice, while the two RL conditions
test whether failed experience and the choice of actor update yield further
benefit.

\begin{table}[t]
  \centering
  \caption{\textbf{Planned real-robot comparison.} Task success rate for a
  shared pretrained FM checkpoint per task and three fine-tuning methods.
  All eight result entries are \emph{not yet measured}; ``---'' is a placeholder,
  not zero success. Evaluation trial counts and uncertainty will accompany
  the completed measurements.}
  \label{tab:real-robot}
  \small
  \setlength{\tabcolsep}{5pt}
  \begin{tabular}{llcc}
    \toprule
    Method & Refinement data & \makecell{Jigglypuff\\into case} &
      \makecell{Red cube\\onto purple cube} \\
    \midrule
    Base FM & None & --- & --- \\
    Success-only BC refinement & Successes only & --- & --- \\
    DICE-RL & Successes and failures & --- & --- \\
    DICE-RL $+$ \method{} (ours) & Successes and failures & --- & --- \\
    \bottomrule
  \end{tabular}
\end{table}

\paragraph{How a robot trial is run.}
The operator places the objects on the specified stair treads, selects a
task/method checkpoint in the evaluation GUI, and resets the robot to the
saved initial configuration. Reset starts each robot from its own measured
joint positions and uses smoothly time-scaled joint-space interpolation at
100\,Hz. Leader-to-follower command forwarding is paused during autonomous
execution, so it cannot compete with policy commands. Model loading and CUDA
warmup finish before task motion begins.

Pressing \emph{Play} starts recording and primes a timestamped action buffer
from live observations. A persistent inference worker generates 30-step
predictions, while a separate control loop sends targets at 15\,Hz. By
default, a fresh inference request is scheduled every eight control periods
($\simeq0.53$\,s); the unused prediction tail provides reserve while inference
runs. Expired targets are discarded, and replacement plans are anchored at
the last command. The current implementation limits arm target changes to
0.6\,rad/s and halts on stale observations, excessive tracking error, or buffer
exhaustion. The GUI shows chest, wrist, and tactile views throughout execution.
These execution settings are shared across the compared policies.

An attempt ends when the operator presses \emph{Stop}, a runtime guard stops
execution, or the common rollout time limit is reached (20\,s in the current
GUI configuration). The operator then labels the attempt \emph{Success} or
\emph{Fail} and chooses \emph{Save} or \emph{Discard}. Before the next attempt,
the operator restores the object arrangement and resets the robot. Valid
failed attempts in the RL collection budget must be saved, with discarded
or invalid attempts accounted for separately to avoid success-selection bias.

\paragraph{Recorded data and refinement.}
Each saved take contains chest/wrist RGB and depth-visualization videos,
separate leader and follower joint-position, Cartesian-pose, and reported
joint-torque logs, executed policy
commands, available tactile measurements, and metadata identifying the
checkpoint, task, termination reason, and binary outcome. Robot trajectory
CSVs are logged at nominally 100\,Hz. Autonomous BC targets and RL replay use
the commands actually issued and follower state, including command limiting
and timing, rather than stationary leader readings or unexecuted predictions.
The success label supplies sparse task-outcome feedback; it is not inferred
from imitation loss or a critic score.

Here \emph{online rollouts} denotes data collected on the physical robot.
The current refinement workflow trains from the saved pool between collection
and evaluation; it does not autonomously acquire new trials during every
gradient-update phase. Thus the real-robot comparison is offline refinement
from online-collected experience. Camera timestamps and bounded histories are
used for live observation alignment. Historical AVI-only recordings lack
per-frame timestamp sidecars and require approximate end-anchored alignment
offline; this limitation will be disclosed for any such data retained in the
controlled study.

\paragraph{Held-out evaluation.}
The planned comparison fixes the starting checkpoint per task, reset and
object-placement distribution, and evaluation conditions. DICE-RL and
DICE-RL $+$ \method{} share the labeled rollout pool, critic architecture,
critic-training procedure, demonstration-replay settings, and optimization
budget and schedule, isolating the actor update. All methods use the same
observation/action interface and execution limits. Fresh evaluation trials
are separate from the base demonstrations, the 40 BC-refinement trajectories,
and the 40 RL collection attempts, and are held out from all training.

We will count Jigglypuff placement as successful when the toy is released
fully inside the case and remains there, and stacking as successful when the
red cube remains supported by the purple cube after release and gripper
withdrawal. A common post-release observation interval and matched trial
counts will be fixed before scoring the four methods. \Cref{tab:real-robot}
reserves the success-rate comparison; the demonstration sequences establish
the task setup but do not supply any of its quantitative entries.
